\documentclass[12pt,a4paper]{report}
\usepackage[]{amsmath,amsthm,amssymb,amscd}
\usepackage[latin1]{inputenc}
%\usepackage[spanish]{babel}
%\usepackage{fullpage}

\topmargin 0pt
\advance \topmargin by -\headheight
\advance \topmargin by -\headsep
     
\textheight 246mm
     
\oddsidemargin 0pt
\evensidemargin \oddsidemargin
\marginparwidth 0.5in
     
\textwidth 159mm


%               Theorem Environments
\theoremstyle{definition}
\newtheorem{ex}{}
\newcommand{\pd}[1]{\frac{\partial}{\partial #1}} %\pd{x}
%               Subquestions numbered with letters
\renewcommand{\theenumi}{\textbf{\alph{enumi}}}

%               Maths definitions

\newcommand{\NN}{\ensuremath{\mathbb N}}
\newcommand{\ZZ}{\ensuremath{\mathbb Z}}
\newcommand{\CC}{\ensuremath{\mathbb C}}
\newcommand{\RR}{\ensuremath{\mathbb R}}
\newcommand{\QQ}{\ensuremath{\mathbb Q}}
\newcommand{\g}{\ensuremath{\frak{g}}}
\newcommand{\h}{\ensuremath{\frak{h}}}
\newcommand{\ft}{\ensuremath{\frak{t}}}
\newcommand{\cB}{\mathcal{B}}
\newcommand{\cN}{\mathcal{N}}
\newcommand{\cL}{\mathcal{L}}
\newcommand{\cD}{\mathcal{D}}
%\newcommand{\pd}[1]{\frac{\partial}{\partial #1}} %\pd{x}
%\newcommand{\pd}[1]{\frac{\partial}{\partial #1}} %\pd{x}

\begin{document}
\noindent
{\sffamily\large Marco Zambon 
\hfill   \\Geometr\'ia III, curso 2011-12}
 

\noindent
\hrulefill{}\\
\begin{center}\bfseries\large\textsf{Hoja 1    \vspace{-2mm}}
\end{center}

\noindent
%\hrulefill{}\\

 

 \begin{ex}Sea $M$ una variedad topol\'ogica de dimension $n$ con borde. Demuestra:
\begin{enumerate}
\item $Int(M)$ es una variedad topol\'ogica de dimension $n$ (sin borde).
\item $\partial M$ es una variedad topol\'ogica de dimension $n-1$ (sin borde).
\end{enumerate}
\end{ex}

\noindent{\bf Soluci\'on:}\\

a)  
Sea $p\in Int(M)$. Por definici\'on de interior,  existe un abierto $U\subset M$ con $p\in U$ y un homeomorfismo $\phi \colon U\to \phi(U)$ a un abierto de $\overline{\RR^n_+}$ tal que $$
\phi(p)\in \RR^n_+:=\{(x_1,\dots,x_n)\in\overline{\RR^n_+}:x_n>0\}.$$
Sea $V$ un entorno abierto  de $\phi(p)$ en $\RR^n_+$. Entonces $$\phi|_{U \cap \phi^{-1}(V)}\colon U \cap \phi^{-1}(V) \to \phi(U)\cap V$$ es un homeomorfismo a 
un entorno abierto de $\phi(p)$ en $\RR^n_+$. Nota que    $U \cap \phi^{-1}(V)\subset Int(M)$ (por definici\'on de interior).

 
Arriba vimos que $\phi(U)\cap V$ es abierto en $\RR^n_+$, y nota que $\RR^n_+$ es abierto en $\RR^n$, as\'i que
$\phi(U)\cap V$ es abierto en $\RR^n$.
En conclusi\'on, hemos demostrado que cada punto de $Int(M)$ tiene un entorno abierto en $Int(M)$ que es homeomorfo a un abierto de $\RR^n$, es decir, que $Int(M)$ es una variedad topol\'ogica de dimension $n$.\\

b) Sea $p\in \partial M$.  Por definici\'on de borde,  existe un abierto $U\subset M$ con $p\in U$ y un homeomorfismo  $\phi \colon U\to \phi(U)$ a un abierto de $\overline{\RR^n_+}$  con $\phi(p)\in \RR^{n-1}\times \{0\}$. Restringiendo el dominio y codominio, obtenemos un homeomorfismo $$\phi|_{\partial M \cap U} \colon \partial M \cap U \to \phi(\partial M \cap U).$$ Nota que $\partial M \cap U$ es un entorno abierto de $p$ en $\partial M$, por la definici\'on de topolog\'ia subespacio. Por el ejercicio 2), para todo $q\in U$ tenemos: $q\in \partial M$ si y solo si $\phi(q)\in \RR^{n-1}\times \{0\}$.
Eso implica que $\phi(\partial M \cap U)=\phi(U)\cap (\RR^{n-1}\times \{0\})$, lo cual -- por definici\'on de topolog\'ia subespacio -- 
es un abierto de $\RR^{n-1}$.

En conclusi\'on, hemos demostrado que $\phi|_{\partial M \cap U}$
es un homeomorfismo de ${\partial M \cap U}$ a un abierto de $\RR^{n-1}$, y por lo tanto que $\partial M$ es una variedad topo\'ogica de dimension $n-1$.
 \end{document}
 %%%%%%%
 \begin{ex}

\end{ex}

 \begin{ex}
\begin{enumerate}
\item 
\end{enumerate}
\end{ex}
 
 
 
 